% Pista alg-23s-q1 · Àlgebra
% Procedència: PAU Catalunya 2023 · Convocatòria de setembre · Sèrie 2 · Qüestió 1
\documentclass[11pt,a4paper]{article}
\usepackage{pista}
\pistainfo{Catalunya 2023}{Convocatòria de setembre}{Sèrie 2}

\begin{document}

\section*{\textcolor{blauPista}{Qüestió 1}}

\noindent Siguin
\[
A = \begin{pmatrix} 2 & 1 \\ 3 & 2 \end{pmatrix}, \quad
B = \begin{pmatrix} 2 & -1 \\ -3 & 2 \end{pmatrix} \quad \text{i la matriu identitat d'ordre dos } I = \begin{pmatrix} 1 & 0 \\ 0 & 1 \end{pmatrix}.
\]

\subsection*{\textit{a)} \normalfont Comproveu que $(A - 2I)^{2} = 3I$. \hfill\textnormal{[0,5~punts]}}

\pista{És un càlcul directe:
\begin{itemize}
\item obtén la matriu $A - 2I$ (resta element a element);
\item eleva-la al quadrat;
\item comprova que el resultat coincideix amb $3I$.
\end{itemize}

\medskip
\textbf{Atenció:} elevar al quadrat una matriu vol dir multiplicar-la per ella mateixa (fila per columna), no elevar al quadrat cada element.}

\subsection*{\textit{b)} \normalfont Utilitzant la igualtat de l'apartat anterior, trobeu la matriu inversa de la matriu $A$ en funció de les matrius $A$ i $I$, i comproveu que coincideix amb la matriu $B$. \hfill\textnormal{[1,25~punts]}}

\pista{La idea clau: si portes la igualtat $(A - 2I)^{2} = 3I$ a la forma
\[
A \cdot M = I,
\]
on $M$ és una expressió en funció d'$A$ i $I$, aquesta $M$ és, per definició, la inversa $A^{-1}$.
\begin{itemize}
\item Desenvolupa el quadrat $(A - 2I)^{2}$ com un binomi: ho pots fer perquè $A$ i $I$ commuten ($A \cdot I = I \cdot A = A$).
\item Deixa a l'esquerra els termes amb $A$ i treu $A$ factor comú, escrivint $4A$ com $A \cdot (4I)$: dins del parèntesi no pots restar un nombre a una matriu.
\item Quan tinguis $M$, calcula-la i comprova que dona $B$.
\end{itemize}

\medskip
\textbf{Atenció:} amb dues matrius qualssevol, $(X + Y)^{2} = X^{2} + XY + YX + Y^{2}$; només si $XY = YX$ queda $X^{2} + 2XY + Y^{2}$.}

\subsection*{\textit{c)} \normalfont Calculeu la matriu $X$ que satisfà la igualtat $A \cdot X = B$. \hfill\textnormal{[0,75~punts]}}

\pista{A l'apartat \textit{b)} ja has vist que $A^{-1} = B$: no la tornis a calcular.
\begin{itemize}
\item Per aïllar $X$ de $A \cdot X = B$, multiplica per $A^{-1}$ a banda i banda, \emph{per l'esquerra}.
\item Quan substitueixis, $X$ es redueix al producte de dues matrius que ja coneixes.
\end{itemize}}

\end{document}
